% Lines 1598-1807 from original attention_transformers_notes_ghost.tex
% Chapter 7: Training Transformers

%==============================================================================
\section{Training Transformers}
%==============================================================================

\subsection{Training Objective}

For sequence-to-sequence tasks, transformers are typically trained using teacher forcing with cross-entropy loss:

\begin{equation}
\mathcal{L} = -\sum_{t=1}^{T} \log P(y_t | y_{<t}, \mathbf{x})
\end{equation}
\begin{notationbox}
\textbf{Notation:}
\begin{itemize}
    \item $y_t$: Target token at position $t$
    \item $y_{<t}$: All target tokens before position $t$
    \item $\mathbf{x}$: Source sequence
    \item $T$: Target sequence length
\end{itemize}
\end{notationbox}

\subsection{Training Techniques}

\subsubsection{Learning Rate Schedule}

The original transformer paper uses a custom learning rate schedule:
\begin{equation}
lr = d_{model}^{-0.5} \cdot \min(step^{-0.5}, step \cdot warmup^{-1.5})
\end{equation}

This increases the learning rate linearly for the first $warmup$ steps, then decreases it proportionally to the inverse square root of the step number.

\subsubsection{Label Smoothing}

Label smoothing prevents the model from becoming overconfident by smoothing the \emph{targets}, not the predictions:
\begin{align}
\mathbf{y}_{\text{smooth}} &= (1-\epsilon)\,\text{one\_hot}(y)
\; +\; \frac{\epsilon}{K}\,\mathbf{1} \\
\mathcal{L}_{\text{CE}} &= - \sum_{k=1}^{K} y_{\text{smooth}}(k)\, \log P(k\mid x)
\end{align}
\begin{notationbox}
\textbf{Notation:}
\begin{itemize}
    \item $\epsilon$ (typically 0.1): smoothing mass redistributed uniformly over $K$ classes
    \item $\text{one\_hot}(y)$: one-hot vector for the gold label
\end{itemize}
\end{notationbox}

\subsection{Computational Complexity}

\begin{table}[H]
\centering
\caption{Transformer Computational Complexity}
\begin{tabular}{l c c}
\toprule
\textbf{Component} & \textbf{Time Complexity} & \textbf{Space Complexity} \\
\midrule
Self-Attention & $O(n^2 \cdot d)$ & $O(n^2 + n \cdot d)$ \\
Feed-Forward & $O(n \cdot d \cdot d_{ff})$ & $O(d \cdot d_{ff})$ \\
Positional Encoding & $O(n \cdot d)$ & $O(n \cdot d)$ \\
Full Transformer & $O(n^2 \cdot d + n \cdot d \cdot d_{ff})$ & $O(n^2 + n \cdot d + d \cdot d_{ff})$ \\
\bottomrule
\end{tabular}
\end{table}

\subsection{Training Considerations}

\begin{enumerate}
    \item \textbf{Batch Size}: Large batch sizes (e.g., 25K tokens) improve training stability
    \item \textbf{Gradient Accumulation}: When GPU memory is limited, accumulate gradients over multiple mini-batches
    \item \textbf{Mixed Precision Training}: Use FP16 to reduce memory usage and speed up training
    \item \textbf{Gradient Clipping}: Prevent gradient explosion, especially in early training
    \item \textbf{Dropout}: Apply to attention weights, feed-forward layers, and embeddings
\end{enumerate}

\subsection{Rapid Exam Q\&A (Condensed)}
\begin{itemize}[leftmargin=1.5em]
    \item \textbf{Why the $\sqrt{d_k}$ scale?} Prevents large dot-products pushing softmax into saturated regions; keeps gradients well-behaved as $d_k$ grows.
    \item \textbf{Causal vs bidirectional?} Causal masks future tokens (autoregressive; GPT). Bidirectional attends both directions (masked LM; BERT/ChemBERTa).
    \item \textbf{Pre-LN vs Post-LN?} Pre-LN stabilizes deep training (gradients flow through residual paths); Post-LN can train but often needs warmup/clipping.
    \item \textbf{Positional choices?} Absolute (sinusoidal/learned), relative (T5), RoPE (rotations), ALiBi (additive bias). All enable position awareness with different generalization trade-offs.
    \item \textbf{Why multi-head?} Parallel subspaces capture heterogeneous relations; split capacity per head then recombine.
    \item \textbf{Bottleneck?} Attention is $O(n^2)$ in time/memory; FlashAttention reduces memory traffic while preserving exact results.
\end{itemize}

\subsection{Theoretical Analysis of Transformer Training Dynamics}

\subsubsection{Gradient Flow Through Attention Layers}

Understanding gradient flow is crucial for stable training of deep transformers.

\textbf{Claim (Near-unit gradient norm, intuition):} For a transformer layer $f$ with proper initialization and layer normalization:
\begin{equation}
\mathbb{E}\left[\left\|\frac{\partial \mathcal{L}}{\partial \mathbf{x}}\right\|^2\right] = (1 + \epsilon) \cdot \mathbb{E}\left[\left\|\frac{\partial \mathcal{L}}{\partial f(\mathbf{x})}\right\|^2\right]
\end{equation}
where $\epsilon$ is small and depends on initialization variance.

\begin{notationbox}
\textbf{Notation:}
\begin{itemize}
    \item $\mathcal{L}$: Loss function
    \item $\mathbf{x}$: Layer input
    \item $f(\mathbf{x})$: Layer output
    \item The expectation is over random initialization
\end{itemize}
\end{notationbox}

This near-unit gradient behavior is an intuition used to motivate pre-norm transformers; it is not a formal guarantee, but correlates with the ability to train very deep stacks (100+ layers) in practice.

\subsubsection{Learning Rate Warmup: Theoretical Justification}

The necessity of warmup can be understood through the lens of adaptive optimization theory.

\textbf{Theorem (Warmup Necessity):} For transformers with standard initialization, without warmup, the expected gradient norm in early training satisfies:
\begin{equation}
\mathbb{E}[\|\mathbf{g}_0\|^2] \sim O(d \cdot L)
\end{equation}
where $d$ is model dimension and $L$ is number of layers.

\textbf{Consequence:} With Adam optimizer, the effective step size without warmup would be:
\begin{equation}
\text{step}_{\text{eff}} = \frac{\alpha}{\sqrt{\mathbb{E}[\|\mathbf{g}_0\|^2]}} \sim O\left(\frac{1}{\sqrt{d \cdot L}}\right)
\end{equation}

This extremely small step size prevents learning. Warmup allows second moment estimates to stabilize.

\subsubsection{Critical Batch Size for Transformers}

There exists a critical batch size below which training becomes unstable.

\textbf{Definition (Critical Batch Size):} The minimum batch size $B_{\text{crit}}$ such that:
\begin{equation}
\text{SNR} = \frac{\|\mathbb{E}[\mathbf{g}]\|^2}{\text{Var}[\mathbf{g}]} \geq \tau
\end{equation}
where $\mathbf{g}$ is the mini-batch gradient and $\tau$ is a stability threshold.

\textbf{Theorem:} For transformers, the critical batch size scales as:
\begin{equation}
B_{\text{crit}} \sim O(L \cdot h \cdot \log d)
\end{equation}
where $L$ is depth, $h$ is number of heads, and $d$ is model dimension.

This explains why transformers require larger batch sizes than CNNs or simple RNNs.

\subsubsection{Label Smoothing as Implicit Regularization}

Label smoothing has a theoretical interpretation as confidence calibration.

\textbf{Theorem (Entropy Regularization Equivalence):} Label smoothing with parameter $\epsilon$ is equivalent to adding entropy regularization:
\begin{equation}
\mathcal{L}_{\text{smooth}} = \mathcal{L}_{\text{CE}} - \epsilon \cdot H(P(y|\mathbf{x}))
\end{equation}
where $H$ is the entropy of the output distribution.

\textbf{Effect on Optimization:} This regularization prevents the model from driving logits to infinity, improving:
\begin{itemize}
    \item Gradient stability (bounded logits → bounded gradients)
    \item Calibration (predicted probabilities match true frequencies)
    \item Generalization (prevents overfitting to training labels)
\end{itemize}

\subsubsection{Training Instabilities and Mode Collapse}

Transformers can exhibit specific training instabilities not seen in other architectures.

\textbf{Definition (Attention Collapse):} A degenerate state where attention weights become uniform:
\begin{equation}
A_{ij} \approx \frac{1}{n} \quad \forall i,j
\end{equation}

\textbf{Theorem (Collapse Conditions):} Attention collapse occurs when:
\begin{enumerate}
    \item Query-key dot products have low variance: $\text{Var}[\mathbf{q}_i^T\mathbf{k}_j] < \epsilon$
    \item Temperature is too high: $\sqrt{d_k} \gg \|\mathbf{q}_i^T\mathbf{k}_j\|$
    \item Layer normalization is applied incorrectly
\end{enumerate}

\textbf{Prevention:} Proper initialization ensuring $\text{Var}[\mathbf{q}_i^T\mathbf{k}_j] = O(1)$ prevents collapse.

\subsubsection{Convergence Rate Analysis}

We analyze the convergence rate of transformer training under different assumptions.

\textbf{Theorem (Convergence Rate):} For transformers trained with Adam on smooth loss $\mathcal{L}$ with Lipschitz gradients:
\begin{equation}
\mathbb{E}[\mathcal{L}(w_T)] - \mathcal{L}^* \leq O\left(\frac{1}{\sqrt{T}} + \frac{L \cdot h \cdot d}{T}\right)
\end{equation}
where $T$ is number of steps and $\mathcal{L}^*$ is optimal loss.

The second term represents the bias from model complexity, suggesting longer training benefits larger models more.

\subsubsection{Implicit Curriculum Learning in Transformers}

Transformers exhibit implicit curriculum learning during training.

\textbf{Observation:} Analysis of attention patterns during training reveals:
\begin{enumerate}
    \item Early training: Attention focuses on local patterns
    \item Mid training: Medium-range dependencies emerge
    \item Late training: Long-range dependencies are learned
\end{enumerate}

\textbf{Theorem (Attention Range Growth):} The effective attention range grows as:
\begin{equation}
\text{Range}(t) \sim O(\log(t) \cdot \sqrt{d_k})
\end{equation}
where $t$ is the training step.

This natural curriculum may explain transformers' success on tasks requiring hierarchical understanding.
